% =============================================================================
% Esqueleto del tercer trabajo de la línea CRAFT.
% Marco de trabajo: formulación bi-objetivo (ingresos del gestor de
% infraestructuras ↔ equidad entre operadores) con factibilidad por
% construcción y marco de evaluación común.
% Compilación: pdflatex main.tex (dos pasadas para las referencias).
% Los resultados de la batería E1-E4 están marcados con \pendiente{...}.
% =============================================================================
\documentclass[11pt,a4paper]{article}

\usepackage[utf8]{inputenc}
\usepackage[T1]{fontenc}
\usepackage[spanish,es-tabla]{babel}
\usepackage{amsmath,amssymb}
\usepackage{booktabs}
\usepackage[margin=2.5cm]{geometry}
\usepackage{xcolor}
\usepackage[colorlinks=true,linkcolor=blue!60!black,citecolor=blue!60!black,urlcolor=blue!60!black]{hyperref}

\newcommand{\pendiente}[1]{\textcolor{red}{[PENDIENTE: #1]}}
\newcommand{\ru}{\textup{RU}}
\newcommand{\im}{\textup{IM}}

\title{De la optimización escalar con asignación heurística a la\\
planificación bi-objetivo con factibilidad por construcción:\\
un marco común para ingresos y equidad en mercados\\
ferroviarios liberalizados\\
\large Borrador de trabajo --- tercer trabajo de la línea CRAFT}

\author{David Muñoz Valero\thanks{Universidad de Castilla-La Mancha.
\texttt{david.munoz6@alu.uclm.es}}}
\date{\today}

\begin{document}
\maketitle

\begin{abstract}
En un mercado ferroviario liberalizado, el gestor de infraestructuras
recibe de las empresas operadoras solicitudes de circulación y debe decidir
qué servicios admite y con qué horarios, con dos objetivos en tensión: los
ingresos por el uso de la infraestructura y la equidad en el reparto entre
operadores. Dos trabajos previos de esta línea abordaron el problema por
separado. El primero maximizaba los ingresos mediante metaheurísticas, con
un flujo de dos etapas: primero se ajustaban los horarios solicitados y
después una regla constructiva decidía qué servicios operaban. El segundo
incorporó la equidad mediante un objetivo único que multiplica los ingresos
por un índice de equidad, fijando así el compromiso entre ambos objetivos
antes de la búsqueda. Este trabajo unifica y corrige ambos enfoques con
tres piezas: (i) una \emph{codificación directa} en la que la decisión de
operar cada servicio y los ajustes de horario evolucionan juntas dentro del
algoritmo genético, con las restricciones estructurales garantizadas por la
propia representación y con los conflictos de separación reevaluados sobre
cada horario propuesto; (ii) una \emph{formulación bi-objetivo} (ingresos
penalizados $\leftrightarrow$ equidad entre operadores) resuelta con un
algoritmo genético multiobjetivo que traza el conjunto completo de
soluciones de compromiso en lugar de fijarlo \emph{a priori}; y (iii) un
\emph{marco de evaluación común} que mide con la misma regla las soluciones
de cualquier método y expone los defectos estructurales del flujo de dos
etapas: conflictos surgidos tras la optimización, servicios operativamente
inviables y equidad que no sobrevive al filtro de operatividad. La batería
experimental abarca 12 instancias (25/50/75 servicios $\times$ 4 escenarios
de cuota entre operadores), 8 métodos, 25 semillas, indicadores
multiobjetivo completos con referencia exacta $\varepsilon$-restricción y
contraste estadístico con corrección por comparaciones múltiples.
\pendiente{Actualizar el resumen con los resultados de la batería y las
conclusiones finales.}
\end{abstract}

\tableofcontents

% =============================================================================
\section{Introducción}
\label{sec:introduccion}

\paragraph{El problema.} En un mercado ferroviario liberalizado conviven
varias empresas operadoras (\emph{railway undertakings}, \ru{}s en la
notación de esta línea) que solicitan al gestor de infraestructuras (IM,
\emph{infrastructure manager}) franjas de circulación sobre una red
compartida. Cada solicitud describe un servicio --- origen, destino,
paradas y horarios deseados --- y el IM debe construir el horario operativo:
decide qué servicios admite y, cuando la capacidad no alcanza o los
horarios chocan entre sí, los ajusta o los descarta. De cada servicio
programado el IM percibe un canon, que se penaliza cuando el horario
concedido se desvía del solicitado (más desvío, menor ingreso); y la
asignación entre operadores debe ser equitativa, un principio que la
normativa europea de espacios ferroviarios recoge expresamente. Son por
tanto dos objetivos concurrentes: los ingresos del IM y la equidad del
reparto.

\paragraph{Los trabajos previos.} Esta línea ha abordado el problema en dos
etapas. El primer trabajo \cite{paper1} formuló la maximización de los
ingresos del IM sobre el simulador robin \cite{robin} y comparó diez
metaheurísticas de la librería mealpy \cite{mealpy} en igualdad de
condiciones; el algoritmo genético resultó el mejor método. Su flujo de
resolución, sin embargo, era de \emph{dos etapas}: la metaheurística solo
ajustaba los tiempos de salida de los servicios solicitados y, terminada la
optimización, una regla constructiva ávida (en inglés, \emph{greedy},
``avaricioso'') decidía qué servicios entraban finalmente en el horario,
descartando de cada pareja de servicios en conflicto el menos rentable.
El segundo trabajo \cite{paper2} incorporó la dimensión de equidad entre
operadores mediante índices clásicos (Jain, Gini, Atkinson) adaptados a
capacidades desiguales, y lo hizo con un objetivo único que multiplica los
ingresos por el índice de equidad: la solución resultante depende así del
peso implícito que esa multiplicación impone al compromiso
ingresos--equidad, elegido \emph{antes} de la búsqueda.

\paragraph{Tres defectos del flujo de dos etapas.} Un marco de evaluación
común --- descrito en la sección~\ref{sec:evaluacion} --- permite verificar
tres problemas estructurales del flujo publicado sobre una instancia de
referencia con 50 solicitudes:

\begin{enumerate}
\item \textbf{Conflictos post-optimización.} En el flujo previo, los
conflictos de separación entre servicios se calculan una sola vez, con los
horarios \emph{solicitados}; durante la optimización de horarios esa matriz
no se actualiza, de modo que mover una salida nunca puede resolver un
conflicto. Las soluciones optimizadas reintroducen conflictos que el modelo
no ve (2,6 en promedio por ejecución).
\item \textbf{Servicios operativamente inviables.} La cota secuencial entre
las salidas intermedias de un tren no se impone al muestrear candidatos:
una parte de los servicios queda con horarios operativamente inviables
(15,4 trenes operativos de media frente a 17 declarados).
\item \textbf{Equidad no operativa.} La equidad que persigue el objetivo
multiplicativo se mide sobre la máscara de servicios \emph{antes} del
filtro de operatividad: la equidad declarada (0,93 en la mejor semilla) se
desmorona al medir sobre los servicios realmente operables (0,40).
\end{enumerate}

\paragraph{El nuevo marco.} Frente a ello, este trabajo propone: (i) poner
la decisión de operar cada servicio y los ajustes de horario \emph{dentro}
del individuo del algoritmo genético (codificación directa), de forma que
las restricciones estructurales se garantizan al decodificar y ningún
candidato inviable consume evaluaciones; (ii) recalcular los conflictos
sobre el horario propuesto en cada evaluación, de modo que mover horarios
sí puede resolver conflictos; (iii) tratar ingresos y equidad como objetivos
explícitos y concurrentes con un algoritmo genético multiobjetivo (NSGA-II
\cite{deb2002}), trazando el conjunto de soluciones no dominadas --- el
\emph{frente de Pareto}: aquellas en las que no se puede mejorar un objetivo
sin empeorar el otro --- en lugar de fijar el compromiso \emph{a priori};
y (iv) medir todas las soluciones, producidas por cualquier método, con el
mismo evaluador.

\paragraph{Contribuciones.} Cuatro contribuciones, numeradas C1--C4; el
resto del texto se refiere a ellas con esas etiquetas:
\begin{description}
\item[C1 --- Codificación directa.] La máscara de servicios y los ajustes
de horario evolucionan juntas; las restricciones estructurales se garantizan
por construcción; los conflictos se reevalúan dinámicamente. Experimento de
ablación (eliminar una pieza del método para aislar su contribución): el GA
del primer paper frente al GA directo con idéntico presupuesto.
\item[C2 --- Reformulación bi-objetivo.] Ingresos penalizados
$\leftrightarrow$ equidad entre \ru{}s con un NSGA-II propio sobre el mismo
genoma y operadores; el frente real frente a la ponderación \emph{a priori}
del segundo trabajo. Ablación: con un solo objetivo, el NSGA-II degenera en
un GA elitista equivalente.
\item[C3 --- Marco de evaluación común.] Verificador de conflictos dinámico
validado uno a uno contra robin, filtro de operatividad e índices de equidad
con capacidades, que convierte los tres defectos anteriores en evidencia
medible.
\item[C4 --- Paquete reproducible.] Familia de 12 instancias versionadas con
manifiesto, herramientas de línea de órdenes y batería de experimentos
reanudable.
\end{description}

% =============================================================================
\section{Modelo del problema y notación}
\label{sec:modelo}

Las \ru{}s solicitan un conjunto $SS=\{S_1,\dots,S_{|SS|}\}$ de servicios
sobre un corredor; cada $S_i$ contiene su ruta y el horario solicitado
$dt_i^j$ (salidas) y $at_i^j$ (llegadas), con $j$ indexando las estaciones
del servicio. El IM decide la máscara de planificación $B_i \in \{0,1\}$ y
los horarios actualizados $\tilde{dt}_i^j$, acotados por:

\begin{itemize}
\item el margen de modificación del IM: $dt_i^1 - \delta \le \tilde{dt}_i^1
\le dt_i^1 + \delta$, con $\delta$ el margen admisible en la salida del
origen (restricción 1);
\item los tiempos de operación: $\tilde{dt}_i^j \ge \tilde{dt}_i^{j-1} +
\tau_i + \sigma_i$, con $\tau_i$ el tiempo de marcha del tramo y $\sigma_i$
el tiempo de parada (restricción 2);
\item las extensiones máximas de parada: $\tilde{dt}_i^j \le dt_i^j +
\Delta_{\max}$ (restricción 3) y $\tilde{dt}_i^j \le \tilde{dt}_i^{j-1} +
\tau_i + \Delta_{\max}$ (restricción 4);
\item los conflictos de separación sobre cada tramo de corredor compartido
en el mismo sentido (restricción 6, sección~\ref{sec:conflictos}), con
$\omega$ la distancia mínima de seguridad entre trenes (\emph{headway}).
\end{itemize}

Los ingresos del servicio $i$ --- el \emph{revenue} del IM en la
terminología de los trabajos previos --- replican el modelo económico del
primer trabajo: el canon $ca_i$ --- la tarifa que el operador pagaría si su
solicitud se programara sin cambios --- se penaliza por el desvío respecto
de lo solicitado según la sensibilidad $k_i$ del operador y los topes
máximos de penalización $P_{DT_i}$ y $P_{TT_i}$, con las curvas
$f(x,k) = 1 - e^{-kx^2}\big(0{,}5\cos(\pi x) + 0{,}5\big)$ aplicadas sobre
el desvío normalizado por el margen $\delta$. La
tabla~\ref{tab:notacion} resume la correspondencia y señala dónde cambia el
tratamiento respecto de los trabajos previos.

\begin{table}[t]
\centering
\small
\begin{tabular}{@{}lll@{}}
\toprule
Símbolo & Concepto & Tratamiento en este marco \\
\midrule
$SS$, $S_i$ & solicitudes de las \ru{}s & instancia sintética versionada \\
$B_i$ & servicio planificado & \textbf{gen del individuo} (antes: regla posterior) \\
$DT$ / $ODT$ & salidas solicitadas / actualizadas & genoma de ajustes $\to$ acumulados \\
$\delta$, $\Delta_{\max}$ & margen del IM, tope de parada & cotas del genoma y del acotado \\
$\tau_i$, $\sigma_i$ & tiempos de marcha y parada & invariantes por construcción \\
$\omega$ & distancia mínima de seguridad & umbral $2\omega$ del conflicto dinámico \\
$ca_i$, $k_i$, $P_{DT}$, $P_{TT}$ & canon y penalizaciones & idénticos al primer trabajo \\
$C_{i,i'}$ & indicador de conflicto & \textbf{dinámico} sobre el horario propuesto \\
$FSS$ & horario final & decodificado del mejor individuo \\
\bottomrule
\end{tabular}
\caption{Correspondencia de notación y dónde cambia el tratamiento. El
vocabulario genético se usa en su sentido estándar: un \emph{individuo} es
una solución candidata completa y su \emph{genoma} su representación
interna; una \emph{generación} es una iteración del algoritmo.}
\label{tab:notacion}
\end{table}

% =============================================================================
\section{Codificación directa y factibilidad por construcción}
\label{sec:codificacion}

Cada individuo es un vector de números enteros (minutos) con un bloque por
servicio, $[B_i, \delta_i^1, \dots, \delta_i^{n_i-1}]$, donde $B_i$ decide
si el servicio se planifica y cada $\delta_i^j$ es el ajuste en minutos
sobre la salida solicitada de la estación con salida $j$. La decodificación
acumula los desplazamientos aguas abajo y los acota (operación
\emph{clip}: recortar al intervalo permitido) de manera que reproduce las
restricciones 1--4 sin ninguna reparación posterior:

\begin{align}
s_i^1 &= \mathrm{clip}\big(\delta_i^1,\; -\delta,\; +\delta\big), \\
s_i^j &= \mathrm{clip}\Big(s_i^{j-1} + \delta_i^j,\;\; s_i^{j-1},\;\;
\min\big(\Delta_{\max},\; s_i^{j-1} + \Delta_{\max} - \sigma_i^j\big)\Big),
\qquad j = 2,\dots,n_i-1, \\
\tilde{dt}_i^j &= dt_i^j + s_i^j, \qquad
\tilde{at}_i^{n_i} = \tilde{dt}_i^{n_i-1} + \tau_i^{n_i-1},
\end{align}

donde $s_i^j$ es el desplazamiento acumulado de la estación $j$. La cota
inferior $s_i^j \ge s_i^{j-1}$ equivale exactamente a exigir
$\delta_i^j \ge 0$: con ajustes no negativos, el tiempo entre salidas
decodificadas es $(\tau_i + \sigma_i) + \delta_i^j \ge \tau_i + \sigma_i$,
que es la restricción 2 escrita en otra coordenada --- el conjunto de
soluciones factibles es idéntico y ninguna evaluación se desperdicia en
horarios malformados. Combinando las cotas 2 y 3 se obtiene además una
restricción derivada, no presente explícitamente en el primer trabajo: en
un servicio con paradas intermedias el origen no puede retrasarse más allá
de $\Delta_{\max}$ (restricción 5), que acota el primer gen del bloque.

La única infactibilidad que un individuo puede arrastrar son los
conflictos, gestionados por penalización directa en el objetivo --- sin
llamar a ninguna heurística que ``arregle'' candidatos dentro del bucle:

\begin{equation}
F(g) \;=\; \sum_{i:\,B_i=1} r_i(g) \;-\;
\lambda \sum_{\substack{(i,i') \in C(g)\\ B_i = B_{i'} = 1}}
\min\big(r_i,\, r_{i'}\big),
\label{ec:fitness}
\end{equation}

donde $r_i(g)$ son los ingresos del servicio con su horario decodificado y
$\lambda = 2$: con $\lambda = 1$, mantener dos servicios en conflicto
rendiría exactamente $\max(r_i, r_{i'})$ --- lo mismo que descartar el más
barato --- y la búsqueda carecería de presión para resolver el conflicto.
La regla constructiva del primer trabajo no se pierde: se conserva como
método de comparación (sección~\ref{sec:metodos}) y se ejecuta fuera del
bucle de optimización.

% =============================================================================
\section{Conflictos dinámicos y su lectura como separación ferroviaria}
\label{sec:conflictos}

Dos servicios comparten tramos del corredor. Para cada tramo compartido
recorrido en el mismo sentido, con $\Delta dt_{i,i'} =
\tilde{dt}_i^{a} - \tilde{dt}_{i'}^{a}$ y $\Delta at_{i,i'} =
\tilde{at}_i^{b} - \tilde{at}_{i'}^{b}$ las diferencias de salida y llegada
en sus extremos --- tiempos interpolados linealmente con la distancia en
las estaciones de paso (hipótesis de movimiento rectilíneo uniforme) --- hay
conflicto cuando:

\begin{equation}
C_{i,i'} = 1 \iff
\Delta dt_{i,i'} \cdot \Delta at_{i,i'} \le 0
\;\lor\; |\Delta dt_{i,i'}| \le 2\omega
\;\lor\; |\Delta at_{i,i'}| \le 2\omega.
\label{ec:conflicto}
\end{equation}

El producto negativo captura los cruces --- un tren sale antes y llega
después, imposible en vía única --- y los umbrales $2\omega$ capturan la
invasión mutua de las zonas de seguridad: cada tren ``lleva'' su margen
$\omega$. La condición se evalúa sobre el \emph{horario propuesto} en cada
evaluación, de modo que desplazar una salida puede eliminar un conflicto
durante la búsqueda; en el flujo previo, con la matriz de conflictos fijada
sobre los horarios solicitados, ese movimiento era invisible. El
verificador propio está validado uno a uno contra la semántica del
simulador robin en la instancia de referencia (103 de 103 pares
idénticos).

\paragraph{Lectura como separación ferroviaria.} La condición $2\omega$
admite una lectura natural como sustituto macroscópico (en inglés,
\emph{surrogate}) de la separación que impone la capa de protección
ferroviaria: la escalera de tiempos de bloqueo (\emph{blocking times}) del
bloque fijo convencional --- cuyo aprovechamiento se estudia con el método
UIC 406 \cite{uic406} ---, las curvas de frenado del sistema europeo de
gestión del tráfico ERTMS en su nivel 2, y el cantón móvil del nivel 3
(\emph{moving block}) \cite{besinovic2021,vanaken2017}. La distancia
$\omega$ parametriza, por tanto, el régimen de protección de la
infraestructura, y el verificador dinámico comprueba la factibilidad del
horario \emph{contra esa capa}: distingue el horario \emph{operable} del
horario ``de papel''. El modelo es deliberadamente macroscópico a nivel de
corredor --- sin cantones, curvas de frenado ni transiciones de nivel --- y
esta conexión se usa solo como encuadre, no como contribución; la
extensión a distancias de seguridad por pareja de trenes
$\omega_{ij} = f(\tau_i, \tau_j)$, la esencia del cantón móvil, queda como
trabajo futuro. El encuadre normativo es el de la Directiva
2012/34/UE \cite{directiva201234}: asignación de capacidad sin
discriminación entre operadores con contraprestaciones por uso.

% =============================================================================
\section{Formulación bi-objetivo: ingresos frente a equidad}
\label{sec:biobjetivo}

El segundo trabajo fijó el compromiso ingresos--equidad \emph{a priori},
multiplicando los ingresos por el índice de equidad elegido de antemano. El
nuevo marco lo trata como lo que es --- un problema con dos objetivos
concurrentes --- con el vector

\begin{equation}
\mathbf{F}(g) = \big(F(g),\; J_{\alpha}\big(B(g)\big)\big),
\qquad J_{\alpha} = \frac{\big(\sum_r \rho_r^{\alpha}\big)^2}
{|\mathcal{R}| \sum_r \rho_r^{2\alpha}},
\end{equation}

donde $J_{\alpha}$ es el índice de Jain calibrado del segundo trabajo
($\alpha = 25$) sobre los ratios sensible a capacidades
$\rho_r = a_r / c_r$ entre la importancia lograda $a_r$ del operador $r$ y
su capacidad $c_r$ (su cuota de servicios solicitados), y $B(g)$ es la
máscara decodificada: la equidad depende solo de \emph{qué} servicios
operan, los ingresos también de \emph{cómo}. El índice de Jain valora un
reparto entre 0 (máxima inequidad) y 1 (equidad perfecta); el exponente de
calificación $\alpha$ penaliza con fuerza los ratios desequilibrados.

El compromiso que el IM extrae del frente se toma en el \emph{punto de
rodilla} del frente (en inglés, \emph{knee}): la solución más próxima a la
esquina ideal tras normalizar ambos objetivos entre 0 y 1 --- el compromiso
en el que ya apenas se gana en un objetivo sin perder mucho en el otro.
La elección del frente frente a la ponderación \emph{a priori} no es
estética: el compromiso aceptable depende del contexto normativo y comercial
de cada corredor, y el frente lo ofrece completo en una sola búsqueda,
mientras que la ponderación exige reoptimizar por cada peso ensayado y solo
garantiza el punto correspondiente a cada peso si la búsqueda escalar fuera
perfecta.

% =============================================================================
\section{Métodos de resolución}
\label{sec:metodos}

Este trabajo propone dos variantes sobre la misma cabeza de problema ---
un algoritmo genético mono-objetivo y un NSGA-II bi-objetivo --- y las
compara con los métodos de los trabajos previos y con referencias externas.
Esta sección describe cada algoritmo, su contexto y su papel.

\subsection{El algoritmo genético con codificación directa}
\label{sec:ga}

El algoritmo genético es el método clásico de optimización poblacional:
una población de soluciones candidatas que evoluciona por selección de los
más aptos, recombinación y variación aleatoria. Sobre la codificación de la
sección~\ref{sec:codificacion} se usan operadores que respetan la
estructura por bloques --- cada bloque contiene la decisión de un servicio
y sus ajustes, y ningún operador los separa:

\begin{itemize}
\item \textbf{Inicialización}: todos los servicios planificados y ajustes
con ruido reducido alrededor de cero, es decir, la población parte de los
horarios solicitados --- el escenario real de planificación y el de
penalización mínima.
\item \textbf{Selección}: torneo de tres individuos (gana el de mayor
ingresos penalizados).
\item \textbf{Cruce}: intercambio de bloques completos de servicio entre
progenitores, de modo que la decisión sobre un servicio viaja siempre con
sus ajustes.
\item \textbf{Mutación}: inversión del bit de planificación (renunciar a un
servicio), perturbación gaussiana entera de un ajuste en minutos, o
devolución del ajuste a cero --- que significa ``dejar el servicio como lo
solicitó el operador'', el movimiento semánticamente más útil.
\item \textbf{Elitismo}: el mejor individuo pasa intacto a la generación
siguiente, lo que garantiza una convergencia que nunca empeora.
\end{itemize}

\subsection{NSGA-II: contexto, funcionamiento y encaje}
\label{sec:nsga2}

\paragraph{Contexto.} La optimización multiobjetivo clásica agregaba los
objetivos en uno solo antes de resolver --- sumas ponderadas o la
restricción $\varepsilon$ --- lo que obliga a fijar o barrer el compromiso
\emph{antes} de conocer el frente. El NSGA-II (\emph{Non-dominated Sorting
Genetic Algorithm II}, Deb et al., 2002 \cite{deb2002}) popularizó el
enfoque evolutivo: una única población busca simultáneamente todas las
soluciones no dominadas y devuelve el frente completo.

\paragraph{Funcionamiento.} Por generación, tras evaluar la población:

\begin{enumerate}
\item \textbf{Ordenación no dominada rápida}: la población se estratifica
en frentes por dominación de Pareto; el primer frente es el conjunto de
soluciones no dominadas, el segundo el que resulta al retirarlo, y así
sucesivamente.
\item \textbf{Distancia de hacinamiento}: dentro de cada frente, la
densidad local de cada solución (distancia media a sus vecinas por
objetivo); los extremos del frente tienen distancia infinita y se
preservan siempre.
\item \textbf{Torneos binarios con comparación jerárquica}: gana la
solución de mejor rango; a igual rango, la menos hacinada --- presión
hacia el frente y hacia su diversidad a la vez.
\item \textbf{Selección ambiental $\mu+\lambda$}: progenitores y
descendientes se fusionan y los supervivientes se toman frente a frente;
el frente que excede el presupuesto se trunca por distancia de
hacinamiento (elitismo: lo mejor nunca se pierde).
\end{enumerate}

\paragraph{Para qué problemas se usa.} Es el algoritmo de referencia de la
optimización evolutiva multiobjetivo y el comparador estándar en la
literatura: diseño de ingeniería con magnitudes en tensión (coste frente a
prestaciones, peso frente a resistencia), planificación y asignación de
recursos, y en general problemas con dos o tres objetivos en conflicto en
los que el decisor quiere el conjunto de compromisos, no un punto. Para más
de tres objetivos se prefieren otras familias de algoritmos; aquí solo hay
dos.

\paragraph{Encaje en nuestro problema.} Cuatro razones concretas: (i) los
ingresos y la equidad están genuinamente en tensión --- admitir servicios
menos rentables de operadores infrarrepresentados mejora la equidad y
reduce los ingresos ---, de modo que existe un frente con contenido;
(ii) el IM no conoce \emph{a priori} el compromiso aceptable, que depende
del contexto normativo y comercial, así que el producto útil es el menú
completo de compromisos; (iii) con dos objetivos, NSGA-II opera en su
régimen más eficaz; y (iv) se integra sobre el mismo genoma y los mismos
operadores que el GA, de modo que ambas variantes difieren únicamente en la
selección y la comparación queda limpia. Con un solo objetivo ($M{=}1$), la
ordenación no dominada colapsa en frentes de un elemento y el NSGA-II
degenera en un GA elitista sobre la población fusionada: esa equivalencia
permite, además, un experimento de ablación del tratamiento multiobjetivo.

\subsection{Por qué se proponen dos variantes}
\label{sec:dosvariantes}

Las dos variantes propuestas responden a las dos preguntas del trabajo y
comparten genoma, operadores y presupuesto, de modo que las diferencias
observables se atribuyen a la formulación y no a la implementación:

\begin{itemize}
\item \textbf{GA directo (mono-objetivo)}: responde a ¿qué aporta la
representación por sí sola? Optimiza los ingresos penalizados de la
ecuación~\eqref{ec:fitness} y se compara con los métodos de dos etapas de
los trabajos previos a igual presupuesto (ablación de la codificación
directa, contribución~C1). Es, además, la continuidad directa con la
formulación del primer trabajo.
\item \textbf{NSGA-II (bi-objetivo)}: responde a ¿qué aporta el tratamiento
multiobjetivo? Añade la equidad como segundo objetivo explícito y se
compara con el barrido escalar del segundo trabajo (ablación de la
formulación bi-objetivo, contribución~C2). La comparación entre $M{=}1$ y $M{=}2$ verifica que el
beneficio proviene de la selección multiobjetivo y no del simple elitismo.
Es, a su vez, la extensión natural de la formulación del segundo trabajo.
\end{itemize}

\subsection{Los métodos de comparación}
\label{sec:comparadores}

Los métodos restantes completan la comparación; cada uno con un papel
definido:

\begin{itemize}
\item \textbf{Regla constructiva ávida del primer trabajo} (en inglés,
\emph{greedy}, ``avaricioso''): parte del horario solicitado completo y
descarta de cada pareja en conflicto el servicio menos rentable,
iterativamente, hasta obtener un horario sin conflictos; no mueve ningún
horario. Representa el techo de la asignación heurística: lo máximo que el
esquema de dos etapas puede conseguir sin optimizar horarios.
\item \textbf{GA de mealpy mono-objetivo}: el GA estándar de la librería
mealpy \cite{mealpy} sobre las salidas reales, con la máscara derivada de
la regla constructiva --- el flujo de dos etapas completo del primer
trabajo, del que era el mejor de los diez algoritmos evaluados. Su
comparación con el GA directo aísla la contribución de la codificación.
\item \textbf{GSA con reparación}: el algoritmo de búsqueda gravitacional
de la línea --- una población de ``masas'' que se atraen según la ley de
la gravitación en proporción a su calidad --- con dos mecanismos añadidos:
los ajustes de cada candidato se convierten en desplazamientos acumulados
factibles (reparación, la misma pieza que la decodificación directa usa por
construcción) y la población inicial se siembra con los horarios
solicitados. Muestra que el segundo algoritmo de la línea, reparado,
alcanza también el techo del esquema de dos etapas: la mejora no depende
del algoritmo concreto sino de la representación.
\item \textbf{Barrido escalar rev$\times$equidad del segundo trabajo}: los
puntos que la formulación anterior obtenía fijando el compromiso \emph{a
priori} (con Jain, Gini y Atkinson y sus exponentes calibrados); son los
puntos que el frente del NSGA-II debe dominar o cubrir.
\item \textbf{NSGA-II de pymoo}, librería externa de referencia en
optimización multiobjetivo: validación independiente de nuestra
implementación --- dos frentes obtenidos con maquinaria distinta sobre el
mismo problema deben superponerse sin dominación mutua masiva (comparada
con la métrica de cobertura, abajo). Otros algoritmos multiobjetivo quedan
fuera del alcance: la tesis no es que NSGA-II sea el mejor algoritmo
multiobjetivo, sino que la formulación bi-objetivo es la correcta.
\item \textbf{Modelo exacto $\varepsilon$-restricción} (Pyomo, lenguaje de
modelado algebraico, con el solucionador SCIP, de programación entera
mixta): impone una equidad mínima $\varepsilon$ como restricción explícita
sobre una formulación exacta de selección y resuelve a optimalidad;
proporciona puntos de referencia óptimos para calibrar los indicadores.
Solo en las instancias de 25 servicios --- más allá su resolución no
escala, lo que se documenta como límite del modelo exacto, replicando el
papel del solucionador exacto en el primer trabajo.
\end{itemize}

% =============================================================================
\section{Marco de evaluación común}
\label{sec:evaluacion}

Toda solución producida por cualquier método --- vector de salidas más
máscara de servicios --- se mide con el mismo evaluador:

\begin{enumerate}
\item \textbf{filtro de operatividad}: solo los servicios cuyos tiempos de
operación decodificados se sostienen ganan ingresos y cuentan en la
asignación por operador;
\item \textbf{conflictos dinámicos} reevaluados sobre el horario propuesto
con el verificador de la sección~\ref{sec:conflictos}, e ingresos penalizados
según la ecuación~\eqref{ec:fitness};
\item \textbf{equidad operativa}: los índices de Jain, Gini y Atkinson con
capacidades y calibración del segundo trabajo, más la inequidad por
diferencias entre pares de operadores, computados sobre la máscara
\emph{operativa} (tras el filtro), no sobre la declarada.
\end{enumerate}

Este marco convierte la observación motivadora en evidencia medible: el
flujo de dos etapas publicado acumula conflictos post-optimización y
servicios inviables, y su equidad declarada difiere de la operativa. La
distinción entre horario operable y horario ``de papel'' es, a la vez, la
lectura como separación ferroviaria del verificador dinámico
(sección~\ref{sec:conflictos}).

% =============================================================================
\section{Diseño experimental}
\label{sec:diseno}

\paragraph{Instancias.} Doce instancias del corredor Madrid--Barcelona,
fruto de cruzar tamaños de 25, 50 y 75 servicios solicitados con 4
escenarios de reparto entre operadores (tres operadores con cuotas
muestreadas 50/25/25 --- la instancia de los trabajos previos --- y cinco
operadores con cuotas exactas equilibradas 20\,\% cada uno, decrecientes
suaves 30/25/20/15/10 y dominantes 55/25/10/5/5). Las cuotas fijan la
capacidad de cada operador: la referencia contra la que se mide la equidad
de su asignación. La familia es \emph{versionada}: cada instancia se genera
una sola vez, con semilla fija, y queda congelada como artefacto acompañado
de un manifiesto con sus cuotas, recuentos reales de servicios por operador
y pares en conflicto solicitados; ningún método remuestrea los datos en
tiempo de ejecución, de modo que todas las comparaciones se hacen sobre
exactamente los mismos datos y cualquier ejecución futura es reproducible.
La instancia de 50 servicios y tres operadores coincide con la de los
trabajos previos, de modo que su línea base queda embebida en la familia.

\paragraph{Métodos.} Los métodos comparados y el porqué de cada uno se
describen en la sección~\ref{sec:metodos}: las dos variantes propuestas
(GA directo y NSGA-II bi-objetivo), los métodos de dos etapas de los
trabajos previos (regla constructiva ávida, GA de mealpy, GSA con
reparación), el barrido escalar rev$\times$equidad del segundo trabajo y
las dos referencias externas (NSGA-II de pymoo y modelo exacto
$\varepsilon$-restricción en Pyomo+SCIP).

\paragraph{Presupuesto y estadística.} Todas las metaheurísticas reciben el
mismo presupuesto --- población de 40 individuos y 100 generaciones, es
decir, 4.000 evaluaciones de la función objetivo --- para que la
comparación dependa solo de la representación y la búsqueda, no del
esfuerzo. Cada celda se ejecuta con 25 semillas y se compara con la prueba
de los rangos con signo de Wilcoxon (para muestras emparejadas por semilla)
con corrección de Holm por familia de comparaciones (el ajuste de los
valores $p$ que evita declarar significativos resultados fruto del azar al
contrastar muchas parejas a la vez), y el tamaño del efecto se cuantifica
con la delta de Cliff (probabilidad de que un método supere a otro,
menos la recíproca, en $[-1, 1]$) con sus bandas de magnitud estándar.

\paragraph{Indicadores multiobjetivo.} Cuatro medidas complementarias
evalúan la calidad de cada frente:

\begin{itemize}
\item \textbf{Hipervolumen}: el área que el frente domina por encima de un
punto de referencia (su esquina peor); mayor es mejor y resume a la vez
convergencia y diversidad.
\item \textbf{Distancia generacional invertida} (IGD, por sus siglas en
inglés, \emph{inverted generational distance}): distancia media de cada
punto de referencia al punto más cercano del frente; menor es mejor y mide
lo cerca que el frente está de la referencia.
\item \textbf{Dispersión $\Delta$ de Deb}: uniformidad del espaciado entre
soluciones contiguas del frente y su alcance hacia los extremos; menor es
mejor.
\item \textbf{Métrica de cobertura} (en inglés, \emph{C-metric}): fracción
de las soluciones de un frente dominadas por algún punto de otro; se usa en
ambas direcciones para comparar frentes.
\end{itemize}

El frente de referencia para la distancia generacional invertida se
consolida como la unión de todos los frentes obtenidos más los puntos del
modelo exacto. Criterio de aceptación de la tesis bi-objetivo: el frente de
NSGA-II domina o cubre lo que el barrido escalar muestrea con todas sus
configuraciones (cobertura alta e hipervolumen del frente al menos igual al
de la unión de puntos escalares).

\paragraph{Resultados preliminares.} La tabla~\ref{tab:baseline} congela la
línea base sobre la instancia de referencia (50 servicios, tres operadores,
103 pares en conflicto, 5 semillas, presupuesto idéntico). Junto a los
ingresos crudos se reportan los ingresos penalizados de la
ecuación~\eqref{ec:fitness} con $\lambda{=}2$: cada pareja de servicios en
conflicto resta el doble del ingreso del más barato de los dos, de modo que
mantener la pareja rinde siempre menos que descartar uno de sus servicios y
la búsqueda tiene presión real para resolver los conflictos. La batería
completa (12 instancias $\times$ 25 semillas) se reporta en
la sección~\ref{sec:resultados}.

\begin{table}[t]
\centering
\small
\setlength{\tabcolsep}{4pt}
\begin{tabular}{@{}lrrrrr@{}}
\toprule
Método & Ingresos & Penalizados & Operativos & Conflictos & Jain oper. \\
 & & ($\lambda{=}2$) & & & \\
\midrule
Regla ávida~\cite{paper1} & 729,53 & 729,53 & 17,0 & 0 & 0,33 \\
GA-mealpy mono~\cite{paper1} & 583,05 $\pm$ 39,54 & 395,44 $\pm$ 117,79 & 15,4 $\pm$ 1,1 & 2,6 $\pm$ 1,8 & 0,37 $\pm$ 0,07 \\
GA-mealpy rev$\times$Jain~\cite{paper2} & 456,30 $\pm$ 67,00 & 339,61 $\pm$ 24,95 & 12,2 $\pm$ 1,9 & 1,6 $\pm$ 1,1 & 0,40 $\pm$ 0,09 \\
GSA con reparación & 728,82 $\pm$ 0,01 & --- & 17,0 $\pm$ 0,0 & 0,0 $\pm$ 0,0 & 0,33 $\pm$ 0,00 \\
\textbf{GA directo} & \textbf{724,68 $\pm$ 27,30} & \textbf{724,68 $\pm$ 27,30} & \textbf{17,6 $\pm$ 0,9} & \textbf{0} & 0,44 $\pm$ 0,15 \\
\textbf{NSGA-II bi-objetivo} & 723,79 $\pm$ 42,74 & 627,34 $\pm$ 35,07 & 17,6 $\pm$ 1,1 & 1,2 $\pm$ 0,8 & \textbf{0,85 $\pm$ 0,16} \\
\bottomrule
\end{tabular}
\caption{Línea base congelada sobre la instancia de referencia (5 semillas,
presupuesto idéntico). NSGA-II reporta el punto de rodilla del frente
(sección~\ref{sec:codificacion}); los ingresos penalizados aplican la
penalización por conflictos de la ecuación~\eqref{ec:fitness} con
$\lambda{=}2$. El GSA con reparación alcanza
el techo del esquema de dos etapas --- la regla ávida sobre los horarios
solicitados --- de modo que el diferencial de la codificación directa está en
la factibilidad sin parches, el número de trenes operativos y la extensión
bi-objetivo.}
\label{tab:baseline}
\end{table}

% =============================================================================
\section{Resultados}
\label{sec:resultados}

\pendiente{Rellenar con la batería completa: (i) la figura principal de cada
instancia (frente de NSGA-II frente a los puntos mono-objetivo y exactos);
(ii) tabla resumen con media $\pm$ desviación, valores $p$ con corrección de
Holm y delta de Cliff; (iii) curvas de hipervolumen y distancia generacional
invertida por generación; (iv) escalabilidad 25/50/75; (v) ablaciones de
codificación (GA con heurística frente a GA directo), de número de objetivos
($M{=}1$ frente a $M{=}2$), del peso de conflicto $\lambda$, de la
dominación por restricciones frente a penalización, del índice de equidad y
su exponente, y del análisis $\omega \times \delta$ interpretado como
regímenes de señalización; (vi) figura de conflictos-en-solicitados frente a
conflictos-en-propuestos y de equidad declarada frente a operativa
(robustez del marco de evaluación común, contribución~C3).}

% =============================================================================
\section{Discusión y limitaciones}
\label{sec:discusion}

El marco fortalece la formulación propia de los trabajos previos: el mensaje
no es que los métodos anteriores sean incorrectos, sino que el evaluador
común --- con verificador dinámico validado y filtro de operatividad ---
revela tres defectos estructurales del flujo de dos etapas y los corrige de
raíz mediante factibilidad por construcción. Limitaciones declaradas: el
modelo es macroscópico a nivel de corredor (la condición de separación es un
sustituto de origen ferroviario, no un modelo de ERTMS); la equidad depende
solo de la máscara de servicios, de modo que el extremo de equidad del
frente puede arrastrar conflictos (se discutirá con la variante de
dominación por restricciones --- tratar los conflictos como restricciones
de dominación en lugar de penalizarlos en el objetivo --- o como
limitación); y el modelo exacto $\varepsilon$-restricción no escala más
allá de 25 servicios, un límite del modelo exacto y no del método propuesto.

% =============================================================================
\section{Conclusiones y trabajo futuro}
\label{sec:conclusiones}

\pendiente{Conclusiones tras la batería. Trabajo futuro ya identificado:
distancias de seguridad por pareja de trenes $\omega_{ij} = f(\tau_i,
\tau_j)$ (esencia del cantón móvil), extensión del modelo exacto,
transferencia a instancias reales de red y estudio de tarifas y cánones
endógenos.}

% =============================================================================
\begin{thebibliography}{9}

\bibitem{paper1}
D.~Muñoz-Valero et al.
\newblock An approach to the timetabling problem in deregulated railway
markets based on metaheuristic algorithms.
\newblock \pendiente{Datos de publicación (revista, año, DOI).}

\bibitem{paper2}
D.~Muñoz-Valero et al.
\newblock Towards equitable rail service allocation through
fairness-oriented timetabling in liberalized markets.
\newblock \pendiente{Datos de publicación (revista, año, DOI).}

\bibitem{deb2002}
K.~Deb, A.~Pratap, S.~Agarwal, T.~Meyarivan.
\newblock A fast and elitist multiobjective genetic algorithm: NSGA-II.
\newblock \emph{IEEE Transactions on Evolutionary Computation},
6(2):182--197, 2002.

\bibitem{mealpy}
T.~Van~Thieu, S.~Mirjalili.
\newblock MEALPY: An open-source Python library for metaheuristic
algorithms.
\newblock \pendiente{Verificar revista y año (IEEE Access?).}

\bibitem{robin}
J.~A.~Martín-Baos et al.
\newblock ROBIN: railway optimization and simulation.
\newblock \pendiente{Cita de software: verificar referencias oficiales del
repositorio.}

\bibitem{besinovic2021}
N.~Bešinović, R.~M.~P.~Goverde.
\newblock Review of railway timetabling developments: planning phases,
goals and future directions.
\newblock \pendiente{Verificar revista y año.}

\bibitem{vanaken2017}
S.~Van~Aken, N.~Bešinović, R.~M.~P.~Goverde.
\newblock Designing alternative railway timetables under infrastructure and
temporal heterogeneity.
\newblock \pendiente{Verificar datos (Transportation Research Part B,
2017?).}

\bibitem{uic406}
UIC.
\newblock Capacity (Leaflet 406) --- método UIC 406 de capacidad basado en
tiempos de bloqueo.
\newblock \pendiente{Verificar edición del folleto que se cita.}

\bibitem{directiva201234}
Parlamento Europeo y Consejo de la UE.
\newblock Directiva 2012/34/UE por la que se establece un espacio único
ferroviario europeo.
\newblock DOUE, 2012.

\bibitem{ertms}
ERA / Comisión Europea.
\newblock ERTMS factsheet: increasing infrastructure capacity.
\newblock \pendiente{Verificar la referencia oficial exacta del documento
(para la cifra $\sim$40\%).}

\end{thebibliography}

\end{document}
